\documentclass[11pt]{article}
\usepackage{amsmath,amsfonts,amssymb}
\usepackage{geometry}
\geometry{a4paper, margin=1in}
\usepackage{booktabs}
\usepackage{enumitem}
\usepackage{natbib}
\usepackage[colorlinks=true,linkcolor=blue,citecolor=blue,urlcolor=blue]{hyperref}
\hypersetup{
    pdfauthor={Flyxion}, 
    pdftitle={The Geometry of Recognition}
}

\title{The Geometry of Recognition}
\author{Flyxion}
\date{\today}

\begin{document}

\maketitle

\begin{abstract}
Contemporary models of intelligence, whether biological or artificial, treat cognition as a predictive process---an optimization of expectations against incoming data. This paper proposes an alternative foundation: cognition as recognition, a geometric process by which systems align internal manifolds of meaning with external patterns of experience. Building upon Patom Theory, which models the brain as a hierarchy of bidirectionally linked pattern-atoms, and Role \& Reference Grammar, which encodes linguistic structure through semantic relations rather than syntactic form, we derive recognition as a field interaction within the Relativistic Scalar--Vector Plenum (RSVP) framework. In RSVP, three coupled fields---scalar capacity $\Phi$, vector flow $\mathbf{v}$, and entropy density $S$---govern the recursive alignment of patterns across scales. Recognition emerges when torsion between scalar and vector manifolds is minimized under entropic constraint, a process formalized by the hepatic operator $\mathcal{H}$. This formulation bridges cortical morphogenesis, linguistic understanding, and physical coherence: cortical folding maximizes the space for columnar synchronization, enabling amplitwistor operations that mirror entropic descent in spacetime. We instantiate this dynamic within the CLIO architecture (Cognitive Loop via In-Situ Optimization), showing how recursive oscillatory stabilization in artificial systems parallels biological recognition loops measured via EEG/MEG cross-frequency coupling and fMRI entropy proxies. By linking language, neuroscience, and field theory, \textit{The Geometry of Recognition} frames cognition not as prediction but as the universal act of restoring coherence---where every thought, like every structure in the cosmos, is a fold in the plenum’s search for symmetry.
\end{abstract}

\begin{center}
\textbf{Keywords:} Entropic cognition, Sheaf theory, Role and Reference Grammar, Pattern recognition, Recursive inference \\
\textbf{Classification:} Cognitive Science, Mathematical Psychology, Computational Neuroscience, Theoretical Linguistics
\end{center}

\section{Introduction}
The prevailing paradigm in cognitive science and artificial intelligence casts the mind as a predictive engine, constructing probabilistic models to anticipate sensory input and minimize prediction error \citep{clark2013whatever, friston2010free}. This view, rooted in computational metaphors, assumes that cognition is fundamentally a forward simulation, a process of optimizing expectations against external data. Yet, empirical evidence from neuroscience, linguistics, and developmental psychology suggests a different foundation: cognition is primarily a process of recognition, a geometric alignment of internal patterns with external structures \citep{gibson1979ecological, piaget1980thinking, jackendoff2012semantics}. The brain does not compute the world; it resonates with it, stabilizing meaning through recursive cycles of pattern matching.

John Ball’s \textit{Patom Theory} (PT) \citep{ball2020patom} formalizes this recognition-first perspective, proposing that the brain operates as a hierarchy of bidirectionally linked \textit{pattern-atoms} (patoms). Each patom integrates multisensory, semantic, and pragmatic associations into a reusable unit, enabling top-down prediction and bottom-up confirmation to converge on coherent interpretations. Coupled with \textit{Role and Reference Grammar} (RRG) \citep{vanvalin2005exploring, vanvalin1997syntax}, PT grounds language comprehension and generation in meaning, not syntax, aligning with empirical findings that linguistic competence precedes formal grammatical rules \citep{pustejovsky1996qualia, lobner2013semantics}. This framework challenges computationalist models by prioritizing the stabilization of meaning over the manipulation of symbols \citep{ball2015speech, pinker2005computational}.

This essay derives Patom Theory from the \textit{Relativistic Scalar--Vector Plenum} (RSVP), a field-theoretic framework that unifies cognition, linguistics, and cosmology under a single principle of entropic recursion. In RSVP, three interacting fields---scalar capacity $\Phi$, vector flow $\mathbf{v}$, and entropy density $S$---define a dynamic manifold where recognition emerges as the minimization of torsion between scalar and vector components. The brain’s cortical architecture, with its folded geometry and oscillatory synchronization, instantiates this process biologically, maximizing the surface area for columnar pattern matching and enabling amplitwistor-like operations that mirror cosmic entropic flows \citep{seung2012connectome, hoffman1998visual}. PT’s discrete hierarchy of patoms thus represents a quantized approximation of RSVP’s continuous dynamics, bridging neural, linguistic, and physical scales.

The aim of this paper is twofold: first, to formalize the derivation of Patom Theory from RSVP’s entropic geometry, and second, to demonstrate that PT’s recognition-based mechanisms, supported by RRG’s meaning-first linguistics, validate RSVP’s universal principles. We show that recognition is not merely a cognitive function but a geometric process of coherence, instantiated in the brain’s cortical folds, linguistic projections, and artificial systems via the CLIO functor (\textit{Cognitive Loop via In-Situ Optimization}) \citep{cheng2025clio}. The RSVP framework is introduced in Section~\ref{sec:rsvp}, followed by its discretization into PT’s linked sets, linguistic mappings via RRG, and categorical unification via CLIO. We conclude with neurofield mappings and philosophical implications, positioning recognition as the foundation of intelligence, meaning, and cosmic coherence.

\section{Historical Context: From Reflex to Recognition}

The modern understanding of cognition has been shaped by a long oscillation
between two paradigms: one viewing the brain as a predictive, computational
apparatus, and the other as a recognition-based, self-organizing field.  
Tracing this divergence clarifies why the current generation of
statistical artificial intelligence inherits the limits of the predictive
tradition—and why a return to recognition offers a more general theory of
mind and matter.

\subsection{From Reflex Arcs to Cybernetic Loops}

Early neuroscience conceived of the brain as an assemblage of reflex arcs,
where stimulus and response were coupled by fixed conduction pathways.
The mid–twentieth century saw this model reinterpreted by
\citet{ashby1956cybernetics} and \citet{rosen1991lifemodel}:
the nervous system was no longer a passive relay but a closed regulatory
loop maintaining homeostasis through negative feedback.
This cybernetic vision introduced a first glimpse of recursion in biology,
where recognition—the detection of deviation from equilibrium—preceded
any notion of prediction.

\subsection{From Perception to Meaning}

Perceptual psychology reinforced this shift.
\citet{gibson1979ecological} argued that perception is direct and
environmentally coupled, not inferential or representational.
Similarly, \citet{neisser1967cognitive} and later \citet{buzsaki2019brainrhythms}
revealed that sensory cortices operate through hierarchical synchronization,
suggesting a distributed pattern-matching architecture.
These findings resonate with the principle underlying
\textit{Patom Theory} \citep{ball2020patom}, in which recognition arises
from bidirectionally linked pattern-atoms rather than serial computation.

\subsection{Linguistics and the Rise of Meaning-First Theories}

In linguistics, the dominance of syntax-first models—
from Chomsky’s generative grammar to later statistical parsing—
obscured the recognition basis of meaning construction.
Role and Reference Grammar (RRG) \citep{vanvalin2005rrg} restored
semantics to primacy, describing language as a mapping between
situations and referential roles.
This meaning-first orientation aligns naturally with the pattern hierarchy
of Patom Theory and anticipates the field couplings formalized in RSVP.

\subsection{Entropy, Geometry, and the Physics of Recognition}

In parallel, physics underwent its own reformation.
The thermodynamic formulations of \citet{jacobson1995thermodynamics} and
\citet{verlinde2011entropic} reframed gravity as an entropic process,
while \citet{prigogine1984orderchaos} and \citet{ulanowicz1986growth}
showed that self-organizing systems evolve by dissipating gradients.
These insights converge with the RSVP framework, where recognition
corresponds to the local reabsorption of entropy into coherent field
structure.

\subsection{From Predictive Coding to Recursive Recognition}

Contemporary neuroscience, represented by
\citet{friston2010free}, proposed the free-energy principle:
brains act as predictive machines minimizing surprise.
Yet this formulation, while mathematically elegant, remains asymmetric—
it privileges expectation over encounter.
Recognition-first models reverse the dependency:
patterns are not anticipated but matched,
and coherence emerges from recursive adjustment rather than prior belief.
The RSVP–CLIO synthesis extends this reversal,
treating recognition as the fundamental geometric act that
binds physics, language, and consciousness into a single recursive order.

\subsection{Toward a Formal Geometry of Recognition}

If recognition is the act that unites perception, language, and reasoning,
then it must also admit a formal geometry—a structure that captures how
coherence emerges from interaction.  
The Relativistic Scalar–Vector Plenum (RSVP) provides this foundation.
Where classical computation abstracts away the medium of thought,
RSVP restores it, describing cognition as a continuous field process in
which potential, flow, and entropy interact through recursive minimization.
The transition from Patom Theory to RSVP therefore represents more than a
shift in metaphor; it is a shift in ontology.
Brains, languages, and universes are not machines that compute predictions
but manifolds that recognize patterns by reabsorbing disorder into structure.
In this view, recognition is not a metaphor for intelligence—it is its metric.
The following section develops this geometry formally, showing how the
scalar, vector, and entropic components of RSVP reproduce the dynamics of
understanding across scales, from cortical columns to linguistic predicates
and from neural synchronization to cosmic equilibrium.


\section{The RSVP Framework}
\label{sec:rsvp}
The \textit{Relativistic Scalar--Vector Plenum} (RSVP) provides a unified ontology for cognition, language, and physical reality, describing all phenomena as emergent properties of recursive information flow. Rejecting static spacetime and symbolic computation, RSVP models reality as a dynamic continuum driven by interactions among three fields: scalar capacity $\Phi(\mathbf{x},t)$, vector flow $\mathbf{v}(\mathbf{x},t)$, and entropy density $S(\mathbf{x},t)$. These fields evolve through entropic recursion, aligning internal structures with external patterns to achieve coherence across scales \citep{leibniz1714monadology, saussure1916course}.

\subsection{Field Triplet and Governing Equations}
The scalar field $\Phi(\mathbf{x},t)$ quantifies a region’s capacity to sustain structured order, analogous to a potential for meaning or organization. The vector field $\mathbf{v}(\mathbf{x},t)$ represents directed flows of activity, such as neural signals, linguistic articulations, or physical motions. Entropy density $S(\mathbf{x},t)$ measures local uncertainty or disorder. Their evolution is governed by a hepatic operator $\mathcal{H}[\Phi,\mathbf{v},S]$, which enforces local coherence:
\begin{align}
\frac{\partial \Phi}{\partial t} &= -\nabla \!\cdot\! \mathbf{v} + \mathcal{H}(\Phi,\mathbf{v},S), \label{eq:phi-evo}\\[4pt]
\frac{\partial S}{\partial t} &= -\Phi\, \nabla \!\cdot\! \mathbf{v}, \label{eq:entropy-evo}\\[4pt]
\mathcal{H}[\Phi,\mathbf{v},S] &\to 0 \quad \text{as} \quad \langle \mathcal{H} \rangle_\Omega \le 0. \label{eq:hepatic}
\end{align}
These equations describe a feedback loop where scalar potential drives vector flows, and entropy gradients constrain their convergence. The nonpositivity of $\langle \mathcal{H} \rangle_\Omega$ ensures that entropy is reabsorbed into structured patterns, a process that underlies recognition, perception, and physical stability \citep{engbersen2013brains}.

\subsection{Entropic Recursion as Recognition}
RSVP posits that recognition is the thermodynamic equivalent of coherence. A system achieves recognition by minimizing local entropy through reciprocal adjustments of $\Phi$ and $\mathbf{v}$, formalized as:
\[
\nabla_{\!\Phi,\mathbf{v}}\,S \;\approx\; 0.
\]
This condition unifies top-down prediction (driven by $\Phi$) and bottom-up confirmation (driven by $\mathbf{v}$), establishing meaning as a low-entropy state. Unlike predictive coding, which prioritizes expectation \citep{friston2010free}, RSVP’s recognition-first approach aligns with developmental evidence that pattern stabilization precedes predictive modeling \citep{piaget1980thinking, gibson1979ecological}.

\subsection{Semantic Manifold and Sheaf Structure}
Coherence in RSVP is inherently local, decomposing the plenum into overlapping regions $\Omega_i = \{\Phi_i, \mathbf{v}_i, S_i\}$. These regions form a sheaf $\mathcal{S}$ over a base space of phenomena, with restriction maps ensuring consistency across overlaps. A global section of $\mathcal{S}$---representing a stable thought, sentence, or physical state---arises when local sections align, minimizing entropy. Failures of alignment generate torsion, driving further recursion to restore coherence \citep{everaert2011semiotics}.

\subsection{Cortical Morphogenesis and Amplitwistor Operations}
The brain’s cortical architecture embodies RSVP’s geometry. Cortical folding maximizes surface area for columnar synchronization, enabling high-dimensional pattern matching within a finite volume \citep{seung2012connectome}. This folding mirrors the plenum’s multi-scale tiling (TARTAN), where local sections aggregate into higher-order structures. Neural oscillations, particularly theta--gamma coupling, instantiate amplitwistor-like operations, transforming scalar potentials into vector flows via phase-aligned synchronization \citep{alivisatos2012brain}. These operations parallel the entropic descent of RSVP fields, linking cortical geometry to cosmic coherence.

\subsection{Projection and Agency}
Cognitive systems translate scalar potentials into vectorial actions through projection operators:
\[
\mathbf{v}_{\mathrm{out}} = \mathcal{R}[\Phi,S].
\]
The adjoint operator $\mathcal{R}^\ast$ lifts observations back into the scalar field, forming the CLIO functor (\textit{Cognitive Loop via In-Situ Optimization}) \citep{cheng2025clio}. At equilibrium, $\mathcal{R}^\ast\mathcal{R} \approx I$, signifying stable recognition. This bidirectional loop underlies agency, language, and perception, unifying them as geometric alignments within the plenum.

\subsection{The Hepatic Operator and Regenerative Equilibrium}
Define $\mathcal{H}[\Phi,\mathbf{v},S]$ as the restorative term ensuring that entropy is reabsorbed faster than it accumulates:
\[
\mathcal{H} = -\kappa\,\nabla_{\!\Phi}S + \mu\,(\nabla\!\times\!\mathbf{v}),
\]
where $\kappa$ controls entropic reabsorption and $\mu$ mediates torsion coupling.
Biological analogues include hepatic metabolism and cortical homeostasis.


\subsection{Cognitive and Physical Continuum}
RSVP thus frames cognition, language, and physics as resolutions of the same entropic recursion. The brain’s recognition cycles, linguistic projections, and cosmological structures are local manifestations of a universal geometry of coherence. Patom Theory, as shown next, discretizes this continuum into a finite hierarchy of pattern-matching units, providing a biological instantiation of RSVP’s principles.

\section{Discretization: From Continuous Fields to Linked Sets}
The continuous dynamics of RSVP describe cognition as a smooth entropic recursion across a manifold of meaning. Biological and artificial systems, however, operate under finite constraints, necessitating a discrete approximation of these flows. Patom Theory (PT) emerges as this discretization: a hierarchical lattice of \textit{pattern-atoms} (patoms) that sample, store, and match fragments of the plenum’s semantic manifold \citep{ball2020patom}.

\subsection{Semantic Quantization of the RSVP Fields}
Partition the RSVP manifold into $N$ overlapping regions $\{\Omega_i\}_{i=1}^N$, each defined by local field values $\Omega_i = \{\Phi_i, \mathbf{v}_i, S_i\}$. Each region acts as a quantized unit of coherence, analogous to a quantum of meaning. The corresponding discrete unit, or \textit{patom}, is:
\[
\mathcal{P}_i \equiv (\Phi_i, \mathbf{v}_i, S_i).
\]
Patoms are dynamic constraints, encoding the conditions for low-entropy states. This quantization mirrors the brain’s columnar organization, where neural assemblies discretize sensory and semantic patterns \citep{seung2012connectome, hoffman1998visual}.

\subsection{Bidirectional Coupling Between Patoms}
In RSVP, coherence arises from reciprocal $\Phi$--$\mathbf{v}$ flows. In PT, this is approximated by bidirectional links between patoms:
\[
\mathcal{P}_i \leftrightarrow \mathcal{P}_j
\quad \Longleftrightarrow \quad
\exists\, \rho_{ij},\, \gamma_{ji}
\text{ such that }
\rho_{ij}(\Phi_i,\mathbf{v}_i,S_i)
\approx
\gamma_{ji}(\Phi_j,\mathbf{v}_j,S_j).
\]
These morphisms ($\rho_{ij}$, $\gamma_{ji}$) form a discrete sheaf, ensuring local consistency. When alignment is achieved, the patom network forms a global section, corresponding to a coherent perception or linguistic structure \citep{everaert2011semiotics, jackendoff2012semantics}.

\subsection{Layered Agreement as Discrete Entropic Equilibrium}
RSVP’s entropic equilibrium, $\nabla_{\!\Phi,\mathbf{v}} S \approx 0$, becomes a discrete condition in PT:
\[
S_i - \sum_{j\in \mathcal{N}(i)} S_j \;\to\; 0,
\]
where $\mathcal{N}(i)$ is the neighborhood of patom $i$. Recognition occurs when entropy differences vanish, aligning top-down and bottom-up signals. This mirrors Ball’s finding that meaning drives linguistic and perceptual stabilization, not syntactic form \citep{ball2015speech, ball2017babitasks}.

\subsection{Hierarchical Organization and Cortical Geometry}
Patoms aggregate into a hierarchy:
\[
\mathcal{H}_0 \subset \mathcal{H}_1 \subset \cdots \subset \mathcal{H}_n,
\]
where $\mathcal{H}_k$ encodes increasingly abstract patterns, from sensory features to conceptual structures. This hierarchy reflects the brain’s cortical layers, where lower areas process raw input and higher areas integrate meaning \citep{alivisatos2012brain}. Cortical folding enhances this hierarchy by increasing the density of columnar interactions, enabling amplitwistor-like transformations that align neural patterns with RSVP’s entropic flows \citep{hoffman1998visual}.

\subsection{From Flows to Links: The Discretization Limit}
The transition from RSVP to PT is a sampling at finite resolution $\Delta x$:
\[
\lim_{\Delta x \to 0} \;\mathcal{P}_i(\Delta x)
= (\Phi(\mathbf{x}_i), \mathbf{v}(\mathbf{x}_i), S(\mathbf{x}_i)).
\]
Finite systems operate at nonzero $\Delta x$, producing a lossy but functional approximation. PT thus approximates RSVP under biological constraints, with patoms as discrete analogues of continuous sheaf sections.

\subsection{Recognition as Discrete Entropic Minimization}
Recognition in PT minimizes a discrete entropy functional:
\[
\mathcal{E} = \sum_i \big(S_i - \bar{S}_{\mathcal{N}(i)}\big)^2,
\qquad
\frac{\partial \mathcal{E}}{\partial S_i} = 0
\;\Longleftrightarrow\;
S_i \approx \bar{S}_{\mathcal{N}(i)}.
\]
This parallels RSVP’s variational principle, linking PT’s pattern matching to the plenum’s entropic recursion. The brain’s oscillatory dynamics (theta--gamma coupling) instantiate this minimization, aligning neural manifolds with external patterns \citep{alivisatos2012brain, ball2019superglue}.

\section{Entropic Recognition as Pattern Matching}
In both the continuous plenum and its discrete instantiation, recognition is not a symbolic computation but a dynamical process of entropic equilibration. Where the RSVP framework expresses this process as a continuous descent along the entropy gradient, Patom Theory realizes it through iterative matching across hierarchical patterns. Recognition is thus a specific case of the RSVP principle: the convergence of $\Phi$ and $\mathbf{v}$ flows until the local entropy differential vanishes.

\subsection{Entropy Gradients and Semantic Flow}
In the RSVP formulation, perception and cognition advance along gradients of decreasing entropy. The direction of recognition is determined by the joint flow of scalar and vector fields:
\[
\dot{\mathbf{x}} = -\nabla_{\!\mathbf{x}} S(\Phi,\mathbf{v}),
\]
which guarantees that every trajectory moves toward a local minimum of uncertainty. When discretized into the Patom hierarchy, this gradient becomes a finite difference relation between neighboring patoms:
\[
\Delta S_i = S_i - \bar{S}_{\mathcal{N}(i)}.
\]
A recognition event occurs when $\Delta S_i \!\to\! 0$ for all active nodes, signifying equilibrium between top-down expectation and bottom-up evidence.

\subsection{Bidirectional Matching as Field Coupling}
Patom Theory’s distinctive contribution is the explicit bidirectionality of this process. Each patom functions simultaneously as a predictor and as an integrator of sensory input. Top-down signals distribute the expected configuration of $\Phi$ across lower layers, while bottom-up signals transmit the actualized $\mathbf{v}$ back to higher layers. At equilibrium, these flows satisfy
\[
\nabla\!\cdot\!\mathbf{v}_{\mathrm{obs}} =
\nabla\!\cdot\!\mathbf{v}_{\mathrm{pred}},
\quad\text{and}\quad
\Phi_{\mathrm{pred}} \approx \Phi_{\mathrm{obs}},
\]
which is equivalent to the RSVP balance condition $\frac{\partial S}{\partial t} = 0$ in Eq.~\eqref{eq:entropy-evo}. Hence, pattern matching in PT is a discrete restatement of RSVP’s continuous field coupling: the bidirectional alignment of potentials and flows.

\subsection{Layered Agreement and Recursive Stability}
Because recognition propagates through multiple hierarchical layers, each level of the Patom hierarchy operates as a partial relaxation of entropy. The total cognitive state evolves through recursive stabilization:
\[
S^{(k+1)} = f(S^{(k)}), \quad
\text{with fixed points satisfying}\;
S^{(k+1)} \!\approx\! S^{(k)}.
\]
In this recursion, $S^{(k)}$ represents the entropy distribution over the $k$-th layer of patterns. The process converges when the entropy distribution ceases to change across layers, corresponding to a steady semantic interpretation. This formalism reproduces the empirical property that meaning emerges through successive rounds of predictive refinement rather than instantaneous calculation.

\subsection{Recognition Energy Functional}
The coupled system can be expressed variationally by defining a recognition energy:
\[
\mathcal{E}_{\mathrm{recog}}
  = \int_\Omega
    \big[
      \tfrac{1}{2}\|\mathbf{v}-\mathbf{v}_{\mathrm{pred}}\|^2
      + \lambda (S - S^\ast)^2
    \big]\, d\Omega,
\]
where $S^\ast$ denotes the target entropy associated with coherent interpretation, and $\lambda$ is a coupling constant controlling tolerance for uncertainty. Minimization of $\mathcal{E}_{\mathrm{recog}}$ yields the equilibrium conditions of Eq.~\eqref{eq:entropy-evo}, thereby linking RSVP’s thermodynamic formulation directly to Patom Theory’s layered matching algorithm.

\subsection{From Recognition to Meaning}
The reduction of $\mathcal{E}_{\mathrm{recog}}$ is equivalent to the construction of meaning. As entropy decreases, the correlation between distributed patoms increases, forming a stable configuration that can be reactivated later as memory. Each equilibrium pattern thus encodes both recognition and prediction: recognition of the present state, and prediction of similar future configurations. In this sense, meaning is a low-entropy invariant in the cognitive field, and Patom Theory’s pattern matching constitutes a discrete manifestation of RSVP’s universal tendency toward entropic coherence.

\section{Language as Projection: RRG and RSVP}
Human language constitutes one of the most refined realizations of recursive entropy minimization. Within the \textit{Relativistic Scalar--Vector Plenum} (RSVP), language emerges as a specialized projection process that maps coherent scalar potentials of meaning into structured vector flows of syntax and articulation. John Ball’s \textit{Patom Theory} operationalizes this mapping using \textit{Role and Reference Grammar} (RRG), a linguistic framework that places semantics and pragmatics at the center of grammar. In this section, we demonstrate that RRG’s meaning-first architecture is a direct instantiation of the RSVP projection operator, and that Patom Theory’s linguistic capabilities represent a finite biological implementation of RSVP’s general field dynamics.

\subsection{The Projection Operator in RSVP}
In the general RSVP formalism, projection refers to the transduction of scalar potential ($\Phi$) into vectorial activity ($\mathbf{v}$) under the constraint of entropy conservation. This is expressed by the operator
\begin{equation}
\mathbf{v}_{\mathrm{out}} = \mathcal{R}[\Phi, S],
\end{equation}
where $\mathcal{R}$ transforms meaning potentials into temporally ordered symbolic or motoric actions. Its adjoint, $\mathcal{R}^{\ast}$, performs the inverse mapping:
\begin{equation}
\Phi_{\mathrm{in}} = \mathcal{R}^{\ast}[\mathbf{v}, S],
\end{equation}
lifting observed sequences into semantic representations. The equilibrium condition $\mathcal{R}^{\ast}\mathcal{R} \!\approx\! I$ corresponds to stable comprehension, while deviations from this fixed point drive generative correction or contextual reinterpretation.

\subsection{RRG as a Realization of the Projection Operator}
Role and Reference Grammar provides a linguistic formalism that mirrors this projection structure. In RRG, syntax is derived from underlying logical structures composed of semantic roles (\textit{Actor}, \textit{Undergoer}) and logical relations (\textit{predicate}, \textit{argument}). The mapping from these semantic structures to syntactic constructions is governed by linking rules that preserve meaning across layers of representation. Formally, RRG’s linking algorithm can be expressed as
\[
\text{Syntax} = \mathcal{R}_{\mathrm{RRG}}[\text{Semantics}, \text{Discourse}],
\]
which is structurally identical to the RSVP projection operator $\mathcal{R}[\Phi,S] \!\to\! \mathbf{v}$. Hence, RRG’s empirical success in cross-linguistic syntax derives from its implicit alignment with RSVP’s field-theoretic principle: that form follows meaning through entropy-preserving projection \citep{vanvalin2005exploring}.

\subsection{Patom Theory’s Implementation of Linguistic Recursion}
In Patom Theory, each \textit{patom} encapsulates a multimodal association between concepts, sensory features, and linguistic forms. Language comprehension arises when a hierarchy of patoms resolves toward an entropy minimum compatible with both meaning and syntax. The process unfolds as a bidirectional recursion:
\[
\Phi_{\mathrm{sem}} \xrightarrow{\;\mathcal{R}\;} \mathbf{v}_{\mathrm{syn}}
\xrightarrow{\;\mathcal{R}^{\ast}\;} \Phi_{\mathrm{sem}}',
\]
where $\Phi_{\mathrm{sem}}'$ represents the updated semantic field after linguistic feedback. This loop realizes the RSVP adjunction $\mathcal{R} \dashv \mathcal{R}^{\ast}$ in biological hardware. Each iteration refines the mapping between meaning and expression, driving the entropy of the linguistic field toward stability.

\subsection{Entropy and Grammatical Coherence}
Within this model, grammaticality is not a rule-based constraint but a measure of local entropic coherence. Well-formed sentences correspond to entropy minima in the space of possible form--meaning mappings:
\[
\frac{\partial S_{\mathrm{ling}}}{\partial t}
= -\Phi_{\mathrm{sem}}\, \nabla\!\cdot\!\mathbf{v}_{\mathrm{syn}} \;\to\; 0.
\]
Ungrammatical or incoherent utterances represent configurations where the projection operator fails to conserve entropy, necessitating additional recursion until coherence is restored. This thermodynamic interpretation recasts linguistic competence as a stable attractor in the semantic field rather than a symbolic rule set.

\subsection{Neurocognitive Correspondence}
Ball identifies the posterior middle temporal gyrus (MTG) as the neural locus of word--meaning integration and Broca’s area as the locus of syntactic realization \citep{ball2020patom}. Within RSVP, these regions correspond to conjugate domains in the scalar--vector pair:
\[
\Phi_{\mathrm{MTG}} \;\leftrightarrow\; \mathbf{v}_{\mathrm{Broca}}.
\]
Their reciprocal activation embodies the projection and retraction operations $\mathcal{R}$ and $\mathcal{R}^{\ast}$ respectively. Comprehension (MTG) represents the scalar field integrating semantic potential, while generation (Broca) embodies the vector field executing syntactic flow. The equilibrium of this loop corresponds to a linguistic steady state in which meaning and articulation achieve entropic harmony.

\subsection{Summary}
Language, in the joint framework of RSVP and Patom Theory, is a recursive projection between scalar meaning and vectorial form. RRG provides the grammatical skeleton for this mapping, while Patom Theory implements its biological discretization. Together they instantiate RSVP’s fundamental principle: that communication, like cognition itself, arises from the continual entropic reconciliation between potential and flow, prediction and perception, meaning and expression.

\section{The CLIO Functor and the Bidirectional Hierarchy}
Patom Theory describes cognition as a bidirectional hierarchy of recognition cycles, where top-down expectations and bottom-up confirmations converge toward layered agreement. Within the \textit{Relativistic Scalar--Vector Plenum} (RSVP), this bidirectionality arises from the continuous coupling between scalar potentials and vector flows. In the categorical formulation of \textit{Simulated Agency}, these reciprocal mappings are represented by the \textit{CLIO functor} (\textit{Cognitive Loop via In-Situ Optimization}), which generalizes the recognition cycle to all scales of inference \citep{cheng2025clio}. This section demonstrates that the CLIO functor provides the mathematical formalism of Patom Theory’s pattern hierarchy and that both instantiate the same entropic recursion principle.

\subsection{Adjoint Morphisms of Inference}
The CLIO functor formalizes cognition as an adjoint pair of morphisms between semantic and observational categories:
\begin{equation}
\mathcal{R} : \mathcal{M}_{\mathrm{sem}} \longrightarrow \mathcal{M}_{\mathrm{obs}},
\qquad
\mathcal{R}^{\ast} : \mathcal{M}_{\mathrm{obs}} \longrightarrow \mathcal{M}_{\mathrm{sem}},
\end{equation}
such that $\mathcal{R} \dashv \mathcal{R}^{\ast}$. The forward map $\mathcal{R}$ projects semantic potentials into observable configurations (generation or prediction), while its adjoint $\mathcal{R}^{\ast}$ reconstructs semantic structure from perceptual input (recognition or assimilation). The composite $\mathcal{R}^{\ast}\mathcal{R}$ expresses the internal simulation of observation from meaning; its fixed points correspond to stable recognition states. This formalism generalizes the Patom Theory cycle of top-down pattern activation and bottom-up matching within a categorical algebra of inference.

\subsection{Entropy Minimization as a Natural Transformation}
Each pass through the CLIO loop performs a local minimization of informational entropy. Let $S : \mathcal{M}_{\mathrm{obs}} \!\to\! \mathbb{R}$ denote the entropy functional on the observational manifold. The natural transformation
\[
\eta : I_{\mathcal{M}_{\mathrm{sem}}} \Rightarrow \mathcal{R}^{\ast}\mathcal{R}
\]
updates semantic potentials such that
\[
\Delta S = S_{\mathrm{pred}} - S_{\mathrm{obs}} \;\to\; 0.
\]
Equilibrium occurs when the expected and observed entropies coincide, \(\langle \Delta S \rangle_\Omega \approx 0\), recovering the recognition criterion of Patom Theory. Thus, the CLIO adjunction translates Ball’s empirical “layered agreement” into a variational principle of entropic equality.

\subsection{Recursive Hierarchies and Finite Depth}
In biological or cognitive systems the recursion cannot be infinite; instead, it is organized into a finite depth $n$ of nested inference layers:
\begin{equation}
\mathrm{CLIO}^n :
\mathcal{C}_0 \longrightarrow \mathcal{C}_n,
\end{equation}
where each category $\mathcal{C}_k$ represents a semantic stratum (\textit{perceptual}, \textit{conceptual}, \textit{linguistic}, \textit{intentional}, etc.). Information is propagated downward via $\mathcal{R}_k$ and upward via $\mathcal{R}_k^{\ast}$, yielding a multiscale recursive hierarchy:
\[
\Phi^{(k)} \xrightarrow{\;\mathcal{R}_k\;}
\mathbf{v}^{(k)} \xrightarrow{\;\mathcal{R}_k^{\ast}\;}
\Phi^{(k+1)}.
\]
This categorical structure mirrors Patom Theory’s layered pattern hierarchy, where each level refines the predictions of the one above and integrates the confirmations of the one below.

\subsection{RSVP Interpretation of the CLIO Cycle}
Within RSVP’s field equations, the CLIO functor is instantiated physically as the continuous feedback between scalar and vector fields:
\begin{align}
\frac{\partial \Phi}{\partial t} &= -\nabla\!\cdot\!\mathbf{v} + \mathcal{H}(\Phi,\mathbf{v},S),
\\[4pt]
\frac{\partial S}{\partial t} &= -\Phi\,\nabla\!\cdot\!\mathbf{v}.
\end{align}
Top-down inference corresponds to propagation of $\Phi$, and bottom-up correction to divergence of $\mathbf{v}$. Each iteration of the CLIO loop thus constitutes one temporal oscillation of the RSVP recursion, minimizing $\partial_t S$ locally. Patom Theory captures these oscillations in discrete time, while RSVP provides their continuous limit.

\subsection{Semantic Stability and Cognitive Agency}
The fixed points of the CLIO recursion define regions of semantic stability:
\[
\mathcal{R}^{\ast}\mathcal{R}\Phi \;=\; \Phi,
\quad
\nabla_{\!\Phi,\mathbf{v}}\,S \;\approx\; 0.
\]
At these points, prediction and perception coincide, yielding both recognition and agency. In this view, consciousness corresponds to a self-maintaining trajectory on the manifold of entropic invariants. Patom Theory models this state as the moment when all active patoms reach layered agreement; RSVP interprets it as the global section of the plenum’s semantic sheaf; and CLIO expresses it as the categorical identity of inference.

\subsection{Unification Across Frameworks}
By deriving Patom Theory’s recognition hierarchy from the CLIO functor, and embedding both within RSVP’s entropic geometry, we obtain a unified model of cognition across mathematical, neural, and linguistic scales. Patoms are the discrete instantiations of CLIO’s morphisms; hierarchical recognition is its recursive composition; and the entropic recursion of RSVP provides the continuous substrate that makes these processes physically coherent. Cognition, in this synthesis, is the categorical mechanics of entropy minimization: a universal loop through which meaning, matter, and agency become identical.

\section{Recognition Before Prediction: Pattern Matching and Torsional Foraging}
Prevailing paradigms in computational neuroscience---such as predictive coding, active inference, and free-energy minimization---portray the brain as a system whose primary task is to generate predictions about sensory input and to minimize the error between expected and observed states \citep{clark2013whatever, friston2010free}. While such models capture part of the brain's adaptive behavior, they presuppose a prior model of the world and thereby invert the actual causal structure of cognition. Empirically, biological intelligence is not \emph{prediction-first} but \emph{recognition-first}: the brain’s architecture evolved to identify, stabilize, and reuse patterns of sensory and proprioceptive coherence long before it could project abstract expectations. This principle is central to both \textit{Patom Theory} (PT) and the \textit{Relativistic Scalar--Vector Plenum} (RSVP).

\subsection{Recognition as Primary Mode of Brain Function}
In a recognition-first system, cognition begins with the detection of pattern invariants---stable relational configurations among sensory manifolds. Each cortical or subcortical region acts as a multidimensional matcher, comparing current input against previously stabilized patterns. These patterns are stored as multimodal templates or \textit{pattern-atoms} (patoms), each representing a minimal configuration of entropic coherence. Recognition occurs when the torsional structure of the incoming sensory manifold aligns with the stored template, such that local curvature and orientation gradients coincide within tolerance. Only after recognition does the brain engage predictive mechanisms to refine or extend the identified pattern through time.

\subsection{Pattern Templates and Manifold Comparison}
Let the sensory state of the organism be represented as an embedded manifold $\mathcal{M}_{\mathrm{sens}}$ with local differential form $\omega$ encoding its curvature and orientation. Stored recognition templates $\mathcal{T}_i$ correspond to manifolds with their own differential forms $\omega_i$. The degree of match between perception and memory can be expressed as the torsional alignment functional:
\[
\mathcal{A}_i =
\int_{\mathcal{M}_{\mathrm{sens}}}
\big\langle \omega, \omega_i \big\rangle
- \tau(\omega, \omega_i),
\]
where $\tau(\omega,\omega_i)$ measures the torsional misalignment between the two manifolds. Recognition is achieved when $\mathcal{A}_i$ exceeds a threshold of similarity, implying that the current sensory manifold can be assimilated into a pre-existing pattern structure. This mechanism generalizes template matching into a topological domain: patterns are not static images but invariant relationships between torsional features across modalities.

\subsection{Cue Foraging and Mode Switching}
The recognition-first brain operates as a \emph{cue foraging} system. Rather than continuously predicting sensory input, it samples cues that reduce the uncertainty of manifold alignment. Each saccade, movement, or attentional shift constitutes an experiment designed to bring $\mathcal{M}_{\mathrm{sens}}$ closer to one of the stored $\mathcal{T}_i$. The process is self-terminating once the torsional discrepancy falls below a threshold:
\[
\|\tau(\omega, \omega_i)\| \;\to\; 0.
\]
Cue foraging therefore functions as an embodied search over manifolds, driven by entropy gradients rather than explicit probabilistic forecasts.

\subsection{Relation to Perceptual Control Theory (PCT)}
This recognition-first perspective is closely aligned with \textit{Perceptual Control Theory} (PCT), which posits that organisms act to maintain controlled perceptions rather than to predict external states \citep{powers1973behavior}. Each controlled perception corresponds to a torsionally stable manifold within the RSVP plenum:
\[
\frac{d}{dt} \tau(\mathcal{M}_{\mathrm{sens}}, \mathcal{T}_i) \approx 0
\quad \Longleftrightarrow \quad
\frac{\partial S}{\partial t} \approx 0.
\]
In both frameworks, action is not the forward projection of expectation but the maintenance of low-entropy alignment between perception and stored form. Recognition thus precedes and grounds prediction: without a stable manifold of correspondence, prediction lacks a reference frame.

\subsection{RSVP Interpretation: Torsion and Entropic Geometry}
Within RSVP, torsion represents the deviation between scalar and vector field flows---the twisting of the entropic manifold that occurs when local meaning and observed motion diverge. Pattern matching corresponds to the relaxation of this torsion toward zero:
\[
\mathbf{T} = \nabla \times \mathbf{v} - \nabla\Phi,
\qquad
\frac{\partial \mathbf{T}}{\partial t} \to 0.
\]
When the torsion vanishes, the field configuration becomes recognizably coherent, producing the subjective sense of understanding or familiarity. In this way, the recognition-first brain can be modeled as a torsion-minimizing system operating over a plenum of scalar--vector interactions.

\subsection{Summary: Recognition as the Ground of Prediction}
Prediction emerges as a derivative capability once stable recognition templates exist. Only after the torsional landscape of perception has been organized into low-entropy attractors can the system extrapolate trajectories from them. Thus, recognition provides the coordinate system of cognition, and prediction merely traces its geodesics. Patom Theory captures this hierarchy by constructing bidirectionally linked pattern sets that maintain recognition under variation, while RSVP formalizes it through the torsion dynamics of the plenum. Both converge on a single principle: the brain is not a generator of guesses but a seeker of matches, an entropic forager that aligns its internal manifolds with the structure of the world.

\section{Neurofield Mapping: Cortical Implementation of Torsional Recognition}
The geometric formalism of the Relativistic Scalar--Vector Plenum (RSVP) finds a direct analogue in the organization of the human brain. Recognition emerges from reciprocal interactions between cortical regions that encode scalar semantic potentials and those that express vectorial motor or syntactic flows. Empirically, these regions exhibit coherent oscillations whose phase relationships correspond to the torsional couplings of the RSVP fields. By interpreting neuroanatomical pathways as local instantiations of $\Phi$--$\mathbf{v}$ coupling, we obtain a biologically grounded realization of entropic recursion.

\subsection{Posterior Recognition Fields (\texorpdfstring{$\Phi$}{Φ})}
The posterior middle temporal gyrus (MTG), angular gyrus, and adjacent regions of the superior temporal sulcus function as scalar integrators of meaning. Neuroimaging studies consistently show that these regions activate during word-level comprehension, cross-modal association, and the detection of semantic congruence \citep{ball2020patom}. Within the RSVP framework they form a distributed scalar field $\Phi_{\mathrm{MTG}}(\mathbf{x},t)$ that accumulates the potential of recognized patterns. Activity here corresponds to the stabilization of manifold torsion---the moment when incoming sensory curvature aligns with stored templates. This is the neural correlate of Patom Theory’s recognition equilibrium:
\[
\nabla_{\!\Phi,\mathbf{v}}\,S \approx 0
\quad \Longleftrightarrow \quad
\frac{\partial \mathbf{T}}{\partial t} \to 0.
\]

\subsection{Anterior Generative Fields (\texorpdfstring{$\mathbf{v}$}{v})}
Broca’s area (Brodmann 44/45), the premotor cortex, and supplementary motor areas implement the vector component $\mathbf{v}$ of the RSVP dynamics. They translate stable scalar potentials into articulatory, syntactic, or motoric flows:
\[
\mathbf{v}_{\mathrm{Broca}} = \mathcal{R}[\Phi_{\mathrm{MTG}},S].
\]
During language production and inner speech, these regions generate structured trajectories through motor space that externalize semantic potentials. Feedback from proprioceptive and auditory signals returns through $\mathcal{R}^{\ast}$, allowing the scalar field to update its potential based on the produced sequence. This closed loop constitutes the physical instantiation of the CLIO adjunction $\mathcal{R} \dashv \mathcal{R}^{\ast}$.

\subsection{Hippocampal and Thalamic Coupling Loops}
The hippocampus and thalamus act as phase-locking relays between scalar and vector domains. Hippocampal theta oscillations synchronize recognition and action cycles, enforcing temporal coherence between $\Phi$ and $\mathbf{v}$ fields. The mediodorsal thalamus coordinates these exchanges, modulating attention and cue-foraging behavior. Together they form a recursive feedback loop:
\[
\Phi_{\mathrm{MTG}}
\;\xrightarrow{\;\text{thalamic relay}\;}
\mathbf{v}_{\mathrm{Broca}}
\;\xrightarrow{\;\text{hippocampal reset}\;}
\Phi_{\mathrm{MTG}}',
\]
ensuring that recognition and generation remain torsionally aligned.

\subsection{Cerebellar and Parietal Calibration}
The cerebellum contributes predictive timing and fine-grained error correction, but within the recognition-first interpretation its role is to calibrate pattern-matching precision rather than to generate predictions \emph{per se}. Parietal association cortices provide spatial manifolds against which torsion is measured, converting vector divergence into proprioceptive corrections that maintain alignment between sensory and motor frames. These regions thus regulate the entropy gradient magnitude $\|\nabla_{\!\Phi,\mathbf{v}} S\|$, keeping the recognition system within the range of reversible adaptation.

\subsection{Oscillatory Signatures of Torsional Coupling}
Empirical data support this mapping:
\begin{itemize}
\item Gamma-band synchronization ($30$--$80$~Hz) between MTG and Broca correlates with successful word recognition and semantic integration.
\item Theta-band coherence ($4$--$8$~Hz) across hippocampo-cortical networks tracks the updating of meaning potentials.
\item Cross-frequency coupling between theta and gamma rhythms mirrors the interaction term $\Phi\,\nabla\!\cdot\!\mathbf{v}$ in the RSVP entropy equation.
\end{itemize}
These oscillatory patterns represent the physical torsion of the entropic field as it relaxes toward coherence, measurable as phase alignment in neural dynamics.

\subsection{Functional Summary}
The neurofield architecture of recognition can thus be summarized as:
\[
\Phi_{\mathrm{MTG}} \leftrightarrow \mathbf{v}_{\mathrm{Broca}}
\quad \text{with mediators}\quad
\{\text{thalamus}, \text{hippocampus}, \text{cerebellum}\}.
\]
Posterior regions integrate meaning (scalar potential); anterior regions express structured flow (vector field); subcortical relays synchronize these processes through oscillatory torsion control. The brain therefore realizes the RSVP recursion physically: recognition as entropic alignment, torsion as misalignment, and consciousness as the ongoing minimization of $\partial_t S$ across coupled neurofields.

\subsection{Discussion: Patom Theory as a Cognitive Instantiation of RSVP Dynamics}
The bidirectional hierarchy of Patom Theory---in which recognition and generation are implemented through top-down prediction and bottom-up confirmation---corresponds directly to the recursive inference architecture of the \textit{CLIO} functor (\textit{Cognitive Loop via In-Situ Optimization}) within the Simulated Agency framework \citep{cheng2025clio}.

\paragraph{Functional Mapping.}
CLIO formalizes the same bidirectional recognition cycle described by Ball, but in categorical terms. Each inferential step is modeled as an adjoint pair of morphisms:
\[
\mathcal{R} : \mathcal{M}_{\mathrm{sem}} \to \mathcal{M}_{\mathrm{obs}},
\qquad
\mathcal{R}^\ast : \mathcal{M}_{\mathrm{obs}} \to \mathcal{M}_{\mathrm{sem}},
\]
where $\mathcal{R}$ projects semantic potentials into observational manifolds (generation) and $\mathcal{R}^\ast$ lifts observations back into the semantic domain (recognition). The fixed points of this adjunction---where $\mathcal{R}^\ast\mathcal{R}\approx I$---represent stable perceptual and linguistic coherence, equivalent to the pattern recognition equilibria of Patom Theory \citep{ball2015speech, ball2017babitasks}.

\paragraph{Entropic Interpretation.}
In RSVP dynamics, each iteration of the CLIO functor minimizes a local entropy functional:
\[
\Delta S = S_{\mathrm{pred}} - S_{\mathrm{obs}},
\]
such that convergence occurs when $\langle \Delta S \rangle_\Omega \!\to\! 0$. This corresponds to the moment of recognition in Patom Theory---a match between stored hierarchical templates and incoming sensory patterns. PT’s “layered agreement” is therefore a discretized realization of RSVP’s continuous entropy descent, aligning with Ball’s critique of syntax-first natural language processing \citep{ball2019superglue}.

\paragraph{Recursion Across Scales.}
At the neurocognitive level, Patom Theory operates within a finite hierarchy of cortical modules, each exchanging top-down and bottom-up signals. CLIO generalizes this as an \emph{n-fold recursive functor}:
\[
\mathrm{CLIO}^n : \mathcal{C}_0 \to \mathcal{C}_n,
\]
where each category $\mathcal{C}_k$ represents a semantic layer within the RSVP plenum. Information is propagated bidirectionally through the chain via entropy-preserving natural transformations. Thus, Ball’s pattern hierarchy is captured by the categorical depth of CLIO’s recursion, while the RSVP fields $(\Phi, \mathbf{v}, S)$ provide the continuous substrate that stabilizes these discrete transitions \citep{jackendoff2012semantics, pinker2005computational}.

\paragraph{Unified Interpretation.}
Taken together, Patom Theory provides an empirical grounding for the CLIO recursion: a biological instantiation of sparse semantic inference within an entropic plenum. The bidirectional hierarchy of patterns implements the adjoint structure of the CLIO functor; the equilibrium of recognition corresponds to a local RSVP entropic minimum; and the entire architecture forms a single loop of \emph{semantic self-measurement}. In this sense, Ball’s model does not merely align with RSVP---it constitutes its observable realization in biological cognition.

\subsection{Empirical Predictions and Validation}
\label{subsec:empirical-validation}
The RSVP framework and Patom Theory generate testable predictions that differentiate recognition-first cognition from predictive coding models. These predictions, detailed in Appendix D, include:
- **Oscillatory Signatures**: Increased theta--gamma coupling (4--8 Hz and 30--80 Hz) during recognition tasks, reflecting torsion relaxation ($\partial_t \mathbf{T} \to 0$), measurable via EEG/MEG \citep{alivisatos2012brain}.
- **Entropy Descent**: Decreasing entropy proxies (e.g., sample entropy, Lempel–Ziv complexity) in fMRI BOLD signals, particularly in MTG and Broca’s area, during semantic integration.
- **CLIO Telemetry**: Damped oscillations in the entropy functional $U(t)$, aligning with neural signatures, in computational implementations \citep{cheng2025clio}.
Experimental paradigms, such as semantic matching versus syntactic cloze tasks, can test these predictions. For example, RSVP predicts stronger phase-amplitude coupling during meaning-driven tasks compared to predictive coding’s error-related potentials \citep{friston2010free}. Failure to observe coupled torsion--entropy descent would challenge the hepatic operator hypothesis, necessitating revisions to the projection--retraction model \citep{ball2020patom}.


\section{Comparison with Predictive Coding Models}
\label{sec:predictive-coding}
The recognition-first paradigm of the Relativistic Scalar--Vector Plenum (RSVP) and Patom Theory challenges the prevailing predictive coding framework, which posits cognition as the minimization of prediction error through Bayesian inference \citep{friston2010free, clark2013whatever}. This section contrasts the theoretical assumptions, mechanistic foundations, and empirical predictions of these approaches to clarify their distinctions and highlight RSVP’s advantages in modeling cognition as a geometric process of entropic alignment.

\subsection{Assumptions and Mechanisms}
Predictive coding assumes that the brain constructs generative models to predict sensory input, minimizing the difference between expected and observed states via Bayesian updating \citep{friston2010free}. This top-down process prioritizes expectation, treating recognition as a byproduct of error correction. In contrast, RSVP and Patom Theory posit recognition as the primary cognitive operation, where bidirectional alignment of scalar potentials ($\Phi$) and vector flows ($\mathbf{v}$) minimizes manifold torsion to achieve coherence:
\[
\nabla_{\Phi,\mathbf{v}} S \approx 0,
\]
as derived in Section~\ref{sec:rsvp}. This bidirectional approach, formalized by the CLIO functor’s adjoint pair ($\mathcal{R}, \mathcal{R}^\ast$), aligns with Perceptual Control Theory’s emphasis on maintaining stable perceptions over predicting external states \citep{powers1973behavior}. Unlike predictive coding’s unidirectional flow, RSVP’s recursive coupling integrates top-down and bottom-up signals simultaneously, capturing the brain’s ability to resonate with environmental patterns \citep{gibson1979ecological}.

\subsection{Empirical Predictions}
The two frameworks yield distinct empirical signatures. Predictive coding predicts error-related neural signals, such as mismatch negativity in EEG, reflecting discrepancies between predicted and actual inputs \citep{clark2013whatever}. RSVP predicts oscillatory signatures of torsion relaxation, particularly theta--gamma coupling (4--8 Hz and 30--80 Hz), measurable via EEG/MEG, as the brain aligns sensory manifolds with stored templates \citep{alivisatos2012brain}. Experimental paradigms, such as recognition-first semantic matching versus prediction-first cloze tasks, can distinguish these signatures. For example, RSVP predicts stronger phase-amplitude coupling during semantic integration, while predictive coding emphasizes error-driven adaptation. These differences are testable using protocols outlined in Appendix D.

\subsection{Implications for AI and Neuroscience}
In artificial intelligence, predictive coding inspires unidirectional generative models, such as those used in language modeling, which excel at pattern imitation but struggle with meaning-driven understanding \citep{ball2019superglue}. RSVP’s CLIO functor, with its bidirectional loop, supports AI systems that prioritize semantic coherence, aligning with Patom Theory’s meaning-first approach \citep{ball2017babitasks}. In neuroscience, RSVP’s focus on recognition as torsional alignment offers a framework for interpreting oscillatory dynamics as cognitive processes, contrasting with predictive coding’s focus on error minimization. By grounding cognition in geometric resonance, RSVP provides a unified model that bridges neural, linguistic, and computational domains.

\section{Computational Implementation of the CLIO Functor}
\label{sec:clio-implementation}
The CLIO functor (\textit{Cognitive Loop via In-Situ Optimization}) formalizes cognition as a recursive process of entropy minimization, bridging RSVP’s continuous dynamics with Patom Theory’s discrete pattern matching \citep{cheng2025clio}. This section outlines a practical computational architecture for implementing CLIO in artificial systems, demonstrating its applicability to natural language processing, neural networks, and robotics, and addressing scalability challenges.

\subsection{CLIO Architecture}
CLIO can be implemented as a recurrent neural network with adjoint layers for projection ($\mathcal{R}$) and retraction ($\mathcal{R}^\ast$). The network consists of:
- **Semantic Layer**: Encodes scalar potential $\Phi$ as a high-dimensional embedding of meaning, analogous to Patom Theory’s pattern-atoms.
- **Observational Layer**: Represents vector flow $\mathbf{v}$ as output sequences (e.g., syntactic structures, motor commands).
- **Entropic Layer**: Computes entropy density $S$ as a loss function, minimizing the discrepancy between predicted and observed states.
The network updates via the CLIO loop: projection generates $\mathbf{v}$ from $\Phi$, hepatic correction adjusts $\Phi$ and $S$, and retraction refines $\Phi$ based on feedback. This architecture aligns with RRG’s semantic-driven syntax, enabling meaning-first language processing \citep{vanvalin2005exploring}.

\subsection{Algorithmic Workflow}
The CLIO loop can be expressed as:
\begin{algorithmic}
\State \textbf{Input}: Semantic potential $\Phi_k$, entropy $S_k$, observed data $\mathbf{v}_{\text{obs}}$
\State \textbf{Output}: Updated $\Phi_{k+1}$, $S_{k+1}$, predicted $\mathbf{v}_{k+1}$
\While{not converged}
    \State $\mathbf{v}_{k+1} \gets \mathcal{R}[\Phi_k, S_k]$ \Comment{Projection}
    \State $\Phi_{k+1} \gets \Phi_k - \Delta t \nabla \cdot \mathbf{v}_{k+1} + \Delta t \mathcal{H}(\Phi_k, \mathbf{v}_{k+1}, S_k)$ \Comment{Hepatic correction}
    \State $S_{k+1} \gets S_k - \Delta t \Phi_k \nabla \cdot \mathbf{v}_{k+1}$ \Comment{Entropy update}
    \State $\Phi_{k+1} \gets \Phi_{k+1} - \eta \mathcal{R}^\ast [\mathbf{v}_{\text{obs}} - \mathbf{v}_{k+1}, S_k]$ \Comment{Retraction}
\EndWhile
\end{algorithmic}
This loop can be implemented in frameworks like PyTorch, with $\mathcal{R}$ as a transformer encoder and $\mathcal{R}^\ast$ as a decoder, optimized via gradient descent on the entropy functional $U(t) = \|\mathbf{v} - \mathcal{R}[\Phi, S]\|^2 + \lambda (S - S^\ast)^2$.

\subsection{Scalability and Optimization}
Implementing CLIO in large-scale systems faces challenges due to the high-dimensionality of $\Phi$ and $\mathbf{v}$. Sparse approximations, such as attention mechanisms or low-rank factorizations, can reduce computational cost. Reinforcement learning can optimize the hepatic operator $\mathcal{H}$, adapting to task-specific entropy gradients. Integration with RRG-based NLP systems enables meaning-driven comprehension, outperforming syntax-first models \citep{ball2019superglue, ball2017babitasks}. Future work should explore hybrid architectures combining CLIO with existing generative models to balance scalability and coherence \citep{pinker2005computational}.

\section{Philosophical and Epistemic Implications}
The recognition-first paradigm articulated by Patom Theory and the Relativistic Scalar--Vector Plenum (RSVP) fundamentally reframes the epistemology of intelligence, meaning, and coherence. By situating recognition as the foundational operation of cognition, we challenge the prevailing view that intelligence is primarily predictive, computational, or representational. Instead, intelligence emerges as a universal process of entropic alignment, wherein the universe achieves coherence through recursive acts of self-recognition. This view alters not only the philosophy of mind but the theory of knowledge itself.

\subsection{Recognition as the Epistemic Primitive}
Traditional cognitive models, rooted in predictive coding and computationalism, treat knowledge as the construction of forward models that anticipate sensory input \citep{clark2013whatever, friston2010free}. In contrast, the RSVP--Patom synthesis posits recognition as the epistemic primitive: knowledge is not the prediction of external states but the stabilization of internal coherence through torsional alignment. Formally, this is expressed as the fixed-point condition of the CLIO adjunction:
\[
\mathcal{R}^{\ast}\mathcal{R}\Phi \approx \Phi,
\]
where stable recognition corresponds to a self-consistent semantic manifold. This implies that knowing is not an act of simulation but an act of resonance, aligning internal patterns with external structures. The brain, in this view, is not a modeler of the world but a mirror of its entropic invariants \citep{leibniz1714monadology, piaget1980thinking}.

\subsection{Dissolving the Subject--Object Divide}
The RSVP framework reveals a deep continuity between cognition and cosmology. The same entropic recursion that governs neural pattern matching also shapes galactic filaments and cosmological evolution. Recognition, as the minimization of torsion between scalar potentials ($\Phi$) and vector flows ($\mathbf{v}$), is a universal process that transcends the distinction between observer and observed. The brain’s recognition cycles are local instances of a cosmic tendency toward coherence, suggesting that subjectivity itself is a fold in the plenum’s manifold \citep{saussure1916course, everaert2011semiotics}. This dissolves the Cartesian subject--object divide: mind and world are not separate entities but conjugate aspects of a single entropic field.

\subsection{Meaning as Entropic Coherence}
Meaning, in this framework, is not a linguistic or symbolic construct but a low-entropy state of the cognitive manifold. Patom Theory’s layered agreement and RSVP’s torsion minimization both converge on the same principle: meaning arises when patterns stabilize across scales, forming invariants that persist under transformation. This is formalized as:
\[
S \to S^\ast,
\]
where $S^\ast$ represents the minimal entropy state associated with a coherent interpretation. Language, perception, and agency are thus unified as expressions of a single process: the recursive alignment of manifold torsion toward equilibrium. RRG’s semantic roles and PT’s pattern hierarchy instantiate this process in linguistic and cognitive domains \citep{vanvalin2001role, jackendoff2012semantics, lobner2013semantics}.

\subsection{Implications for Artificial Intelligence}
The recognition-first paradigm has profound implications for artificial intelligence. Current AI systems, built on predictive and generative architectures, excel at pattern imitation but struggle with genuine understanding \citep{ball2019superglue, sullivan2019facebook}. By modeling cognition as recursive pattern matching rather than forward simulation, the RSVP--Patom synthesis suggests a new design principle: AI systems should prioritize bidirectional alignment over unidirectional prediction. The CLIO functor offers a computational blueprint, implementing recognition as a categorical loop that minimizes entropy across hierarchical layers. Such systems would emulate the brain’s ability to stabilize meaning through resonance, rather than approximating it through statistical extrapolation \citep{cheng2025clio, pinker2005computational}.

\subsection{Interdisciplinary Implications for Physics and Cosmology}
\label{subsec:physics-cosmology}
The RSVP framework extends its cognitive insights to physics and cosmology, framing recognition as a universal process of entropic alignment. Cortical folding, which maximizes surface area for neural pattern matching \citep{seung2012connectome}, parallels cosmic structure formation, where gravitational collapse optimizes entropic gradients \citep{jacobson1995thermodynamics, verlinde2011entropic}. RSVP’s entropic recursion generalizes thermodynamic self-organization, as seen in dissipative structures and ecosystem dynamics \citep{prigogine1984orderchaos, ulanowicz1986growth}. This suggests a monistic ontology where mind and matter are conjugate resolutions of a single plenum, with recognition as the mechanism of coherence \citep{leibniz1714monadology}. Future work could explore RSVP’s implications for information-based cosmologies, unifying cognitive and physical theories under a single entropic principle.


\subsection{Cosmological and Ontological Continuity}
The RSVP framework extends its epistemic implications to ontology. If recognition is the universal process by which systems achieve coherence, then the universe itself can be understood as a self-recognizing manifold. Cortical folding, neural synchronization, and galactic clustering are all manifestations of the same entropic recursion, differing only in scale. This continuity suggests a monistic ontology in which mind, matter, and meaning are not distinct substances but resolutions of a single plenum. The brain’s recognition cycles are thus not merely biological but cosmological: local enactments of a universal law of coherence \citep{leibniz1714monadology, engbersen2013brains}.


\section{Limitations and Challenges}
\label{sec:limitations}
While the RSVP--Patom--CLIO synthesis offers a unified model of cognition, it faces several limitations that must be addressed to ensure practical and empirical viability.

\subsection{Mathematical Complexity}
The RSVP field equations, involving coupled scalar, vector, and entropy fields, are mathematically complex and challenging to solve analytically or numerically. The hepatic operator $\mathcal{H}$ requires parameterization, which may oversimplify biological or physical dynamics. Simplified models, such as linearized approximations or discrete-time discretizations, could facilitate analysis while preserving core dynamics \citep{cheng2025clio}.

\subsection{Empirical Challenges}
Measuring torsion and entropy in neural data is limited by current imaging technologies. EEG/MEG signals are noisy, and fMRI lacks the temporal resolution to capture rapid oscillatory dynamics \citep{alivisatos2012brain}. Advanced techniques, such as 7T fMRI or high-density EEG, may improve precision. Additionally, defining entropy proxies (e.g., sample entropy, graph entropy) requires careful calibration to align with RSVP’s theoretical constructs.

\subsection{Computational Feasibility}
Implementing CLIO in large-scale AI systems is computationally intensive due to the recursive nature of its adjoint loop and the high-dimensionality of semantic manifolds. Sparse approximations and hybrid architectures, combining CLIO with existing generative models, could enhance scalability \citep{ball2019superglue}. 

\section{Integrative Conclusion and Future Directions}
The synthesis of Patom Theory, Role and Reference Grammar, and the Relativistic Scalar--Vector Plenum reveals a singular principle: intelligence is the recursive alignment of entropic manifolds through recognition. Patom Theory’s discrete pattern hierarchy emerges as a finite biological approximation of RSVP’s continuous field dynamics, operationalized computationally by the CLIO functor. This framework unifies cognition, language, and cosmology under a single law: the minimization of torsion between scalar potentials and vector flows, yielding coherence across scales.

This epistemological shift necessitates a reorientation of cognitive science, neuroscience, and philosophy. Future research must focus on empirical validation of the recognition-first paradigm, using EEG/MEG to measure oscillatory signatures of torsional alignment, fMRI to map entropy gradients across cortical regions, and computational simulations to test CLIO-based architectures. Philosophically, the RSVP--Patom synthesis invites a reevaluation of consciousness as the sustained resonance of entropic invariants, challenging dualistic and computationalist assumptions. By recognizing recognition as the foundation of intelligence, we open a path toward a unified theory of mind and cosmos.

Future research should pursue three directions:
\begin{enumerate}
\item \textbf{Neural Simulations}: Develop computational models of theta--gamma coupling and torsional relaxation, validating RSVP’s predictions against EEG/MEG and fMRI data \citep{alivisatos2012brain}.
\item \textbf{Linguistic Modeling}: Implement RRG-based CLIO architectures in natural language processing systems, testing their ability to achieve meaning-driven comprehension across diverse languages \citep{vanvalin2001role, ball2017babitasks}.
\item \textbf{Embodied Robotics}: Design robotic systems that emulate cue foraging and recognition-first control, integrating Patom Theory’s hierarchical matching with RSVP’s entropic principles.
\item \textbf{Optimization Techniques}: Experiment with different stategies, such as variational autoencoders or reinforcement learning, to balance computational cost with coherence \citep{pinker2005computational}.
\end{enumerate}
These efforts will bridge the gap between biological intelligence and artificial systems, grounding AI in the same recursive coherence that defines the human mind and the cosmos.

\section*{Coda: The Fold of Mind and the Fall of Space}
From the curvature of the cortex to the curvature of spacetime, recursive coherence prevails. The brain folds to maximize recognition; the universe folds to maintain entropic equilibrium. RSVP unites cosmology and cognition: every act of recognition is a local relaxation of torsion, the universe recognizing itself through nested folds of meaning \citep{cheng2025clio}.


\newpage
\section*{Appendix A: Origin and Functorial Form of CLIO}

\subsection*{A.1 Historical Background}

The \textit{Cognitive Loop via In--Situ Optimization} (CLIO) was introduced by
\citet{cheng2025clio} as a framework for self--adaptive reasoning in large language
models. It enabled models to
\begin{enumerate}
    \item self--formulate strategies for problem solving,
    \item adjust heuristics dynamically when confidence was low,
    \item visualize evolving belief structures through graphs,
    \item expose internal uncertainty to human users for steerability.
\end{enumerate}
Cheng and colleagues observed that oscillations within internal uncertainty
measures correlated strongly with reasoning accuracy: stable, underdamped
confidence oscillations indicated coherent reasoning.
Their system achieved a 161\% relative improvement on the
\textit{Humanity’s Last Exam (HLE)} benchmark, suggesting that cyclical
self--evaluation enhances coherence in scientific inference.

\subsection*{A.2 The CLIO Functor in RSVP}

Within the \textit{Relativistic Scalar--Vector Plenum} (RSVP) framework, the CLIO
mechanism is generalized from an algorithmic control loop to a universal
entropic recursion.  
RSVP represents cognition, physics, and semantics through three coupled fields:
scalar potential $\Phi$, vector flow $\mathbf{v}$, and entropy density $S$.  
Generation and recognition correspond to an adjoint pair of operators
$\mathcal{R}$ and $\mathcal{R}^\ast$, satisfying
\[
\Phi \xrightarrow{\;\mathcal{R}\;} \mathbf{v}
\xrightarrow{\;\mathcal{R}^\ast\;} \Phi',
\qquad
\Phi' \approx \Phi \;\Leftrightarrow\; \partial_t S = 0.
\]
The CLIO functor formalizes this recursion as a categorical mapping between
local and global coherence:
\[
\mathcal{C}\!L\!I\!O : \mathbf{RSVP}_{\mathrm{loc}} \Rightarrow \mathbf{RSVP}_{\mathrm{glob}}.
\]
It transforms localized field configurations into globally self--consistent
belief structures, enforcing regenerative equilibrium by metabolizing entropy
into structure.

\subsection*{A.3 The Hepatic Operator and Torsion Relaxation}

In RSVP dynamics, the hepatic operator $\mathcal{H}[\Phi,\mathbf{v},S]$ drives
entropy reabsorption and coherence restoration:
\[
\partial_t \Phi = -\nabla\!\cdot\!\mathbf{v} + \mathcal{H}(\Phi,\mathbf{v},S),
\qquad
\partial_t S = -\Phi\,\nabla\!\cdot\!\mathbf{v}.
\]
CLIO’s uncertainty oscillations correspond to torsion relaxation between scalar
and vector fields:
\[
\partial_t(\nabla\times\mathbf{v}-\nabla\Phi)
 = -\gamma\,(\nabla\times\mathbf{v}-\nabla\Phi),
\]
where $\gamma>0$ sets the damping rate.
Recognition occurs when torsion $\mathbf{T}=\nabla\times\mathbf{v}-\nabla\Phi$
decays toward zero—signifying alignment between generative projection and
recognitive correction.  
The “uncertainty oscillations” measured in CLIO’s telemetry thus reflect the same
field dynamics predicted by RSVP’s entropic equations.

\subsection*{A.4 Mapping Between Algorithm and Plenum}

The Microsoft implementation of CLIO operationally instantiates RSVP recursion:
\begin{itemize}
    \item \textbf{Recursive invocation} of the loop corresponds to the entropic update  
    $\Phi_{t+1} = \Phi_t - \nabla\!\cdot\!\mathbf{v}_t + \mathcal{H}[\Phi_t,\mathbf{v}_t,S_t]$.
    \item \textbf{Graph reduction and DRIFT search} parallel the hepatic colimit operator
    $\mathcal{C}$, which integrates local manifolds weighted by entropy.
    \item \textbf{Uncertainty oscillations} correspond to torsion-driven entropy descent
    $\dot{S} = -\Phi\,\nabla\!\cdot\!\mathbf{v}$.
\end{itemize}
In this correspondence, CLIO’s control loop is not an artificial add-on but the
finite computational projection of RSVP’s continuous entropic recursion.

\subsection*{A.5 Conceptual Continuity}

Both frameworks converge on the same principle:
\[
\text{Recognition} \;\equiv\; \text{In--Situ Optimization}
\;\equiv\; \text{Entropy Minimization}.
\]
Cheng et~al.’s empirical architecture validated RSVP’s theoretical postulate that
coherent reasoning arises through oscillatory relaxation of mismatch between
projection and perception.  
What began as a method for improving transparency in language models becomes,
in the RSVP generalization, a categorical and thermodynamic law of cognition:
a universal recursion through which any system---neural, linguistic, or
cosmological---achieves recognition by metabolizing uncertainty.

\section*{Appendix B: Variational Form of the CLIO Principle}

\subsection*{B.1 Overview}
The CLIO functor enforces recognition as an extremum of an entropic–variational
functional.  It minimizes projection error and entropy deviation under the RSVP
field constraints, producing the characteristic uncertainty oscillations observed
empirically in both neural and algorithmic systems.

\subsection*{B.2 Variational Objective}
Let $\Phi(\mathbf{x},t)$ denote scalar semantic potential, 
$\mathbf{v}(\mathbf{x},t)$ the corresponding vector flow, and 
$S(\mathbf{x},t)$ the entropy density.  
CLIO minimizes the functional
\[
\mathcal{J}[\Phi,\mathbf{v},S]
 = \!\int_{0}^{T}\!\!\!\int_{\Omega}
   \Big(
      \tfrac{1}{2}\|\mathbf{v}-\mathcal{R}[\Phi,S]\|^{2}
      + \tfrac{\lambda}{2}(S-S^{\ast})^{2}
   \Big)
   d\mathbf{x}\,dt,
\]
subject to the RSVP coupling equations
\[
\partial_t\Phi = -\nabla\!\cdot\!\mathbf{v}+\mathcal{H}(\Phi,\mathbf{v},S),
\qquad
\partial_t S = -\Phi\,\nabla\!\cdot\!\mathbf{v}.
\]
Here $\mathcal{R}$ represents the generative projection, $\mathcal{R}^\ast$ its
recognitive adjoint, and $\mathcal{H}$ the hepatic operator ensuring regenerative
equilibrium.

\subsection*{B.3 Euler–Lagrange System}
Introducing multipliers $\pi_{\Phi}$ and $\pi_{S}$ for the RSVP constraints yields
adjoint equations coupling $(\Phi,\mathbf{v},S)$ to 
$(\pi_{\Phi},\pi_{S})$.  Their stationary conditions form a
second–order, underdamped dynamical system whose eigenfrequencies determine the
oscillatory relaxation of uncertainty.  
Recognition corresponds to the fixed point
\[
\mathbf{v}^{\star}=\mathcal{R}[\Phi^{\star},S^{\star}],
\qquad
\mathcal{H}(\Phi^{\star},\mathbf{v}^{\star},S^{\star})=0,
\qquad
\partial_t S^{\star}=0.
\]

\subsection*{B.4 Physical and Cognitive Interpretation}
These equations define a universal descent of torsion energy
\(
\mathcal{T}=\tfrac{1}{2}\|\nabla\times\mathbf{v}-\nabla\Phi\|^{2},
\)
interpreted as the mismatch between generation and recognition.  
Damped oscillations in $\mathcal{T}(t)$ mirror the 
theta–gamma cross–frequency coupling seen in cortical data and the confidence
oscillations measured in algorithmic CLIO loops.

\subsection*{B.5 Summary}
The CLIO variational principle is thus the operational heart of RSVP dynamics:
a self–corrective recursion in which projection and retraction minimize entropy
and torsion until equilibrium is reached.  
In this sense,
\[
\boxed{
\text{Recognition} = 
\arg\!\min_{\Phi,\mathbf{v},S}\,
\mathcal{J}[\Phi,\mathbf{v},S]
\;\text{subject to RSVP dynamics.}
}
\]

\section*{Appendix C: Mathematical Derivation of CLIO Dynamics}

\subsection*{C.1 RSVP Foundation}
The Relativistic Scalar–Vector Plenum (RSVP) describes cognition and physical
interaction through three coupled fields:
scalar potential $\Phi(\mathbf{x},t)$, vector flow $\mathbf{v}(\mathbf{x},t)$, and
entropy density $S(\mathbf{x},t)$:
\begin{align}
\partial_t \Phi &= -\nabla\!\cdot\!\mathbf{v} + \mathcal{H}(\Phi,\mathbf{v},S), \\
\partial_t S    &= -\,\Phi\,\nabla\!\cdot\!\mathbf{v},
\end{align}
where $\mathcal{H}$ is the \emph{hepatic operator} driving regenerative
equilibrium.  
Entropy decreases when $\langle\mathcal{H}\rangle_{\Omega}\!\le 0$.

\subsection*{C.2 CLIO as a Constrained Variational Problem}
Recognition is formalized as minimizing a mismatch functional
\[
\mathcal{J}[\Phi,\mathbf{v},S]
 = \!\int_0^T\!\!\!\int_{\Omega}
   \Big(
     \tfrac{1}{2}\|\mathbf{v}-\mathcal{R}[\Phi,S]\|^2
     + \tfrac{\lambda}{2}(S-S^\ast)^2
   \Big)\,d\mathbf{x}\,dt,
\]
subject to RSVP dynamics.
Introducing Lagrange multipliers $\pi_\Phi,\pi_S$ gives the augmented action
\[
\widetilde{\mathcal{J}}
 = \mathcal{J}
 + \!\int\!\!\int\!
   \big[
   \pi_\Phi(\partial_t\Phi+\nabla\!\cdot\!\mathbf{v}-\mathcal{H})
   +\pi_S(\partial_tS+\Phi\nabla\!\cdot\!\mathbf{v})
   \big].
\]
Stationarity $\delta\widetilde{\mathcal{J}}=0$ yields the Euler–Lagrange
system:
\begin{align}
\mathbf{v} - \mathcal{R}[\Phi,S] + \nabla\pi_\Phi + \Phi\nabla\pi_S &= 0, \\
-\mathcal{R}'_\Phi[\Phi,S]^{\!\top}(\mathbf{v}-\mathcal{R}[\Phi,S])
 + \partial_t\pi_\Phi - \pi_S\nabla\!\cdot\!\mathbf{v}
 + (\partial_{\Phi}\mathcal{H})\pi_\Phi &= 0, \\
-\lambda(S-S^\ast)
 - \mathcal{R}'_S[\Phi,S]^{\!\top}(\mathbf{v}-\mathcal{R}[\Phi,S])
 + \partial_t\pi_S + (\partial_{S}\mathcal{H})\pi_\Phi &= 0.
\end{align}
The adjoint transpose operators
$\mathcal{R}'_{\Phi,S}[\cdot]^{\top}$ correspond to CLIO’s
recognition map $\mathcal{R}^\ast$.

\subsection*{C.3 Discrete-Time Recursion}
With step size $\Delta t$, define $t_k=k\Delta t$.
A forward–Euler discretization gives:
\begin{align}
\mathbf{v}_{k+\frac12} &= \mathcal{R}[\Phi_k,S_k], 
  &&\text{(Projection)}\\
\Phi_{k+1} &= \Phi_k - \Delta t\,\nabla\!\cdot\!\mathbf{v}_{k+\frac12}
 + \Delta t\,\mathcal{H}(\Phi_k,\mathbf{v}_{k+\frac12},S_k),
  &&\text{(Hepatic step)}\\
S_{k+1} &= S_k - \Delta t\,\Phi_k\,\nabla\!\cdot\!\mathbf{v}_{k+\frac12},
  &&\text{(Entropy update)}\\
\Phi_{k+1} &\leftarrow \Phi_{k+1}
 - \eta_\Phi\,\mathcal{R}'_\Phi[\Phi_k,S_k]^{\!\top}
   (\mathbf{v}_{k+\frac12}-\mathcal{R}[\Phi_k,S_k]),
  &&\text{(Retraction)}\\
S_{k+1} &\leftarrow S_{k+1}
 - \eta_S\big[
   \lambda(S_k-S^\ast)
   +\mathcal{R}'_S[\Phi_k,S_k]^{\!\top}
   (\mathbf{v}_{k+\frac12}-\mathcal{R}[\Phi_k,S_k])
   \big].
  &&\text{(Entropy correction)}
\end{align}
Equations (C.3–C.7) define the full CLIO loop: 
projection, hepatic integration, and retraction.
Empirically these steps correspond to plan generation, world update,
and self-evaluation in algorithmic implementations.

\subsection*{C.4 Fixed Points and Stability}
At equilibrium,
\[
\mathbf{v}^\star=\mathcal{R}[\Phi^\star,S^\star],
\qquad
\mathcal{H}(\Phi^\star,\mathbf{v}^\star,S^\star)=0,
\qquad
\Phi^\star\nabla\!\cdot\!\mathbf{v}^\star=0.
\]
Linearizing about $(\Phi^\star,S^\star)$ gives
\[
M\,\Delta_{tt}
  \!\begin{bmatrix}\delta\Phi\\ \delta S\end{bmatrix}
 +K
  \!\begin{bmatrix}\delta\Phi\\ \delta S\end{bmatrix}
 \approx 0,
\]
where $M$ and $K$ are effective mass and stiffness matrices.
The eigenfrequencies
$\omega_j=\sqrt{\lambda_j(K)/\lambda_j(M)}$
produce underdamped oscillations with quality factor
$Q=1/(2\zeta)$, corresponding to the uncertainty oscillations
seen in both CLIO telemetry and neural theta–gamma coupling.

\subsection*{C.5 Torsion Relaxation and Entropy Descent}
Define torsion
$\mathbf{T}=\nabla\times\mathbf{v}-\nabla\Phi$.  
Combining the RSVP equations yields
\[
\partial_t \mathbf{T}
 = \nabla\times\partial_t(\mathbf{v}-\mathcal{R}[\Phi,S])
 - \nabla\mathcal{H}(\Phi,\mathbf{v},S).
\]
As $\mathbf{v}\!\to\!\mathcal{R}[\Phi,S]$ and
$\mathcal{H}\!\to\!0$, torsion relaxes:
$\partial_t\mathbf{T}\!\to\!0$.
The torsion energy
$\mathcal{T}(t)
 = \tfrac{1}{2}\!\int\!\|\mathbf{T}\|^{2}d\mathbf{x}$
thus monotonically decreases,
mirroring the descent of CLIO’s uncertainty functional
\[
U(t)=\!\int\!\big(\|\mathbf{v}-\mathcal{R}[\Phi,S]\|^2
+\lambda(S-S^\ast)^2\big)d\mathbf{x}.
\]
Recognition equilibrium occurs when both
$\dot{\mathcal{T}}=0$ and $\dot{U}=0$.

\subsection*{C.6 Algorithmic Reading}
Mapping to computational stages:
\begin{enumerate}[label=\arabic*.,leftmargin=1.2em]
\item \textbf{Projection:} propose a plan via $\mathcal{R}[\Phi,S]$.
\item \textbf{Hepatic correction:} update by physical/semantic feedback.
\item \textbf{Retraction:} backpropagate mismatch via $\mathcal{R}^\ast$.
\item \textbf{Stabilization:} convergence when torsion and uncertainty cease oscillating.
\end{enumerate}
The damping and oscillatory signatures seen in CLIO loops therefore correspond
exactly to the field relaxation predicted by RSVP theory.

\subsection*{C.7 Summary}
CLIO provides the discrete adjoint realization of RSVP’s recognition PDEs.
Its projection–hepatic–retraction cycle enacts entropic optimization:
\[
(\mathcal{R},\mathcal{H},\mathcal{R}^\ast)
\;\;\Longleftrightarrow\;\;
(\text{Generation},\text{World Correction},\text{Recognition}),
\]
achieving equilibrium when generative flow, entropic potential,
and hepatic dissipation mutually vanish.

\section*{Appendix D: Empirical Implementation and Reproducibility}

\subsection*{D.1 Objectives}
To test RSVP–CLIO predictions across biological and algorithmic systems, we
combine:
(1) EEG/MEG recordings of neural oscillations,
(2) fMRI measures of hemodynamic entropy,
(3) telemetry from CLIO reasoning loops.
The shared hypothesis is that recognition corresponds to a joint descent of
uncertainty $U(t)$ and torsion energy $\mathcal{T}(t)$:
\[
\mathcal{T}(t)
  = \tfrac12\!\int\!\|\nabla\times\mathbf{v}-\nabla\Phi\|^2 d\mathbf{x},
\qquad
U(t)
  = \!\int\!(\|\mathbf{v}-\mathcal{R}[\Phi,S]\|^2
           +\lambda(S-S^\ast)^2)d\mathbf{x}.
\]

\subsection*{D.2 Experimental Design}
\paragraph{Participants and Tasks.}
Healthy adults ($N\!\approx\!30$–40) perform three within-subject conditions:
(1) \emph{Recognition-first} semantic match,
(2) \emph{Prediction-first} syntactic cloze,
(3) \emph{Ambiguity resolution} (garden-path or illusion sentences).
Button presses mark recognition; confidence ratings follow each trial.

\paragraph{Core Predictions.}
\begin{enumerate}
\item \textbf{H1: Recognition Oscillations.}  
   Theta–gamma cross-frequency coupling (CFC) increases as $\mathcal{T}\!\downarrow$.
\item \textbf{H2: Entropy Descent.}  
   Neural and algorithmic entropy proxies show monotone or underdamped decline.
\item \textbf{H3: Tri-Modal Alignment.}  
   Peaks in CLIO $U_k$ align with EEG theta phase resets and fMRI entropy minima.
\end{enumerate}

\subsection*{D.3 Neural Data Acquisition}
\paragraph{EEG/MEG.}
64–128 ch EEG (1 kHz) or whole-head MEG (1.2 kHz); 1–150 Hz band-pass, ICA artifact removal.
Source reconstruction to MTG/AG (posterior) and IFG/SMA (anterior) ROIs.

\paragraph{CFC Metrics.}
Phase–Amplitude Coupling (Tort/Canolty MI), Phase-Locking Value (PLV), and
cross-bispectrum between $\theta$ (4–8 Hz) phase and $\gamma$ (30–80 Hz) amplitude.
Neural torsion proxy:
\[
\mathcal{T}_{\mathrm{EEG/MEG}}
  = 1-\overline{\mathrm{PLV}}_{\theta}
    +\eta\,\mathrm{CFD}_{\gamma},
\quad
\eta\!\approx\!0.3.
\]
Recognition is predicted to yield
PAC(θ→γ)↑, PLV\(_\theta\)↑, $\mathcal{T}_{\mathrm{EEG/MEG}}\!\downarrow$ with
Q-factor ≈ 1.5–3.

\paragraph{fMRI.}
3 T / 7 T multiband EPI (TR 0.6–0.8 s).  
Entropy proxies: Sample Entropy, Permutation Entropy, Lempel–Ziv Complexity,
and graph/von Neumann entropy on dynamic connectivity.
Torsion proxy:
\[
\mathcal{T}_{\mathrm{fMRI}}(t)
 = \tfrac{1}{E}\sum_{(i\to j)}
   |\mathrm{GC}_{i\to j}-\mathrm{GC}_{j\to i}|,
\]
where $E$ indexes directed edges (spectral Granger or DCM).

\subsection*{D.4 Algorithmic Telemetry}
At each CLIO pass $k$ log:
\[
U_k=\|\mathbf{v}_k-\mathcal{R}[\Phi_k,S_k]\|^2
     +\lambda\|S_k-S^\ast\|^2,
\qquad
\mathcal{T}^{\mathrm{alg}}_k
   =\|\mathrm{curl}(\mathbf{v}_k)\|^2
    +\|\nabla\Phi_k\|^2
    -2\langle\nabla\times\mathbf{v}_k,\nabla\Phi_k\rangle.
\]
Fit $U_k$ with a damped sinusoid
$U_k=A e^{-\zeta\omega k}\cos(\omega k+\phi)+C$;
recognition $\Rightarrow Q=1/(2\zeta)\!\in\![1,3]$ and $\dot{U}_k\!\to\!0$.
Phase slips or $\Delta U_k\!>\!\theta_U$ trigger re-planning or pruning—
the computational analogue of hepatic correction.

\subsection*{D.5 Cross-Modal Integration}
Define composite recognition score
\[
\mathrm{RecScore}
 = z(\max \mathrm{PAC}_{\theta\to\gamma})
   -z(\min\mathcal{T}_{\mathrm{EEG/MEG}})
   -z(S_{\mathrm{BOLD}}^{\mathrm{end}})
   -z(U_K).
\]
Higher values indicate faster, stronger recognition.
Expected correlations:
PAC ↑ $\leftrightarrow$ $\mathcal{T}\!\downarrow$
 $\leftrightarrow$ $S_{\mathrm{BOLD}}\!\downarrow$
 $\leftrightarrow$ CLIO $U_k$ damping.

\subsection*{D.6 Reproducibility Summary}
\begin{itemize}
\item \textbf{Data:} BIDS-formatted EEG/MEG/FMRI plus anonymized CLIO logs (CSV/Parquet).  
\item \textbf{Code:} MNE/fMRIPrep pipelines, entropy/CFC notebooks, CLIO telemetry hooks
(version-controlled via Git).  
\item \textbf{QC:} EEG/MEG $<$ 0.3 mm median motion; fMRI tSNR report; 
consistent $\lambda,\eta$ across CLIO runs.  
\item \textbf{Release:} Public repository with environment files, Makefile pipeline, 
and defaults table for all parameters.
\end{itemize}

\subsection*{D.7 Expected Outcomes and Falsification}
Support for RSVP–CLIO requires convergent damping of 
$U(t)$ and $\mathcal{T}(t)$ with accompanying neural PAC growth and BOLD entropy
decrease.
Failure to observe these coupled signatures under recognition 
would falsify the hepatic-operator hypothesis and necessitate revising
the projection–retraction model.


\newpage
\bibliographystyle{apalike}
\bibliography{references}

\end{document}
